\chapter{Conclusions and Future Work}
\label{ch:conclusion_futurework}

\section{Conclusion}
\label{ch:conclusion}

本研究針對二維矽光子元件之反向設計問題，
建立一套結合 Factorization Machine (FM) 與 Simulated Annealing (SA) 的
Factorization Machine with Simulated Annealing (FMSA) 最佳化架構，
並進一步導入 Gaussian Upsampling、
2D Fourier Basis Encoding 與
$\beta$-Variational Autoencoder ($\beta$-VAE)
三種結構編碼方法，
探討不同結構表示方式對設計空間維度、
代理模型學習能力與最終光學性能之影響。

研究結果顯示，
結構編碼方式對 FMSA 最佳化表現具有顯著影響。
在彎曲波導設計中，
Gaussian Upsampling 與 2D Fourier Basis Encoding
皆能在有限設計變數下有效產生具有連續導波特性的結構，
其中 Gaussian Upsampling 的最高穿透率可達約 95\%，
而 2D Fourier Basis Encoding 可進一步達到約 99\%。
此結果顯示，
適當的結構編碼可在控制最佳化維度的同時，
保留足夠的幾何表示能力，
進而提升代理模型與退火搜尋的最佳化效率。

在交叉波導設計中，
本研究透過重新定義品質因數(Figure of Merit, FoM)
並加入四重對稱幾何限制，
成功設計出可支援 TE$_0$、TE$_1$ 與 TE$_2$
三種模態傳輸的緊湊型交叉波導。
結果顯示，
Gaussian Upsampling 搭配適當的幾何限制，
能夠有效描述具有高度對稱性與連續導波路徑的結構，
並在有限元件尺寸下兼顧傳輸效率與串擾抑制。
相較於相關文獻，本研究所獲得之交叉波導設計在元件尺寸、
傳輸效率與串擾表現方面具有良好競爭力，
顯示所建立之 FMSA 反向設計流程具有應用於多模態矽光子元件設計的潛力。

另一方面，
不同編碼方式亦存在明顯的適用限制。
在模態分波解多工器設計中，
Gaussian Upsampling 所產生之結構主要具有平滑且連續的幾何特徵，
較難同時形成大尺度導波路徑與局部細微結構，
因此在需要複雜模態散射與耦合機制的設計問題中受到限制。
此結果顯示，
具有良好幾何連續性的編碼方式，
未必適合所有類型的矽光子反向設計問題。

在 $\beta$-VAE 策略方面，
雖然生成式模型能將高維結構映射至較低維潛在空間，
但本研究採用之 FMSA 架構仍需在離散二元變數上建立 QUBO 模型，
使潛在空間受到二元化限制。
當潛在變數數量增加時，
FM 所需學習之設計空間亦快速擴大，
在固定資料量與模擬退火採樣數下，
代理模型難以充分學習設計變數與光學性能之間的關係，
因此限制了 $\beta$-VAE 策略的最佳化表現。
此結果亦顯示，
降低原始結構維度並不必然能提升 FMSA 最佳化效率，
潛在空間的表示形式與維度仍需配合代理模型與退火演算法之特性進行設計。

綜合上述結果，
本研究發現 FMSA 應用於二維矽光子反向設計時，
其最佳化能力並非單純由設計變數數量決定，
而是受到
\textbf{設計空間維度、結構表示能力、訓練資料量以及代理模型學習能力}
等因素共同影響。
過度增加設計變數雖可提升結構自由度，
但亦會增加 Factorization Machine 學習與退火搜尋的困難；
相反地，過度限制結構表示方式則可能無法描述元件所需的局部幾何特徵。
因此，結構編碼的核心在於
\textbf{降低最佳化複雜度與保留結構表示能力之間取得平衡}。

本研究並未將所有編碼策略完整應用於每一種矽光子元件，
而是依據不同結構表示方式的特性，
選擇具代表性的設計問題進行驗證。
彎曲波導主要用於比較不同編碼方法與設計維度之影響；
交叉波導則用於驗證 Gaussian Upsampling
結合幾何對稱條件下之最佳化能力；
模態分波解多工器則進一步呈現平滑型結構編碼
在多尺度幾何需求下的限制。
因此，本研究的目的並非單純比較三種編碼方法的性能高低，
而是藉由不同元件的設計結果，
建立結構表示方式與矽光子設計需求之間的關係。

綜合而言，本研究之主要貢獻可歸納如下：

\begin{enumerate}

\item \textbf{建立可應用於二維矽光子元件的 FMSA 反向設計架構。}

本研究將 Factorization Machine 與 Simulated Annealing
應用於二維矽光子結構設計，
並利用電磁模擬結果建立代理模型，
形成一套不依賴梯度資訊的結構最佳化流程。

\item \textbf{建立結構編碼方式與 FMSA 最佳化能力之間的關係。}

透過 Gaussian Upsampling、
2D Fourier Basis Encoding 與 $\beta$-VAE 三種表示方式，
本研究顯示不同編碼方法在降低設計維度、
保留幾何自由度以及代理模型學習難度方面具有不同特性，
並指出有效的結構編碼必須在
\textbf{最佳化複雜度與結構表示能力}之間取得平衡。

\item \textbf{建立不同結構表示方法於矽光子元件設計中的適用性分析。}

透過彎曲波導、交叉波導與模態分波解多工器之設計結果，
本研究顯示平滑型編碼適合描述連續導波路徑，
而具有較高結構自由度的表示方式則較適合需要局部細節之設計問題。
研究結果可作為未來使用代理模型與退火演算法進行
二維矽光子反向設計時，
選擇結構表示方法與設計維度之參考。

\end{enumerate}
\section{Future Work}
\label{ch:futurework}

基於此研究之設計方法與成果，未來可以朝以下幾個方向持續發展。

首先，在生成式結構表示方面，
本研究結果顯示 $\beta$-VAE 雖具備將高維結構映射至低維潛在空間的能力，
但其潛在表示方式仍受到 FMSA 所需二元 QUBO 形式的限制。
未來可進一步研究更適合離散最佳化的生成式潛在表示方法，
例如可直接產生離散潛在變數之生成模型，
或重新設計代理模型與最佳化架構，
使生成模型的潛在空間能夠被更有效地探索。

其次，不同編碼方式於結構表示能力上具有互補特性。
Gaussian Upsampling 擅長描述平滑且連續的整體幾何輪廓，
而 2D Fourier Basis Encoding 則能提供較高的局部結構自由度。
因此，未來可進一步發展 Hybrid Encoding，
以不同表示方式分別描述全域幾何與局部細節，
在控制最佳化維度的同時，
提升複雜矽光子結構之表示能力。

此外，未來可建立更完整的編碼方法評估基準，
將 Gaussian Upsampling、
2D Fourier Basis Encoding 與 $\beta$-VAE
應用於相同元件，
並統一 FoM、設計區域、設計變數數量、
訓練資料量與最佳化計算預算，
以更系統性地分析不同結構表示方法
在光學性能、最佳化效率與可擴展性方面之差異。

最後，本研究目前主要以二維FDTD模擬驗證所提出之方法。
未來可進一步將 FMSA 反向設計架構應用於
模式轉換器、偏振分光器、波長分波器等其他光子積體電路元件，
並延伸至三維電磁模型。
此外，可將最小特徵尺寸、製程誤差與結構容忍度等限制
直接納入最佳化流程，
並透過元件製作與實際量測驗證模擬結果，
進一步評估所提出方法於實際矽光子設計與製造中的應用能力。
